% =====================================================================
%  Osvaldo Fabian Ramiro Balboa — Representación del espacio de búsqueda y ACO
%  Compilar: latexmk -pdf investigacion.tex
% =====================================================================
\documentclass[11pt,letterpaper]{article}

\usepackage[utf8]{inputenc}
\usepackage[T1]{fontenc}
\usepackage[spanish,es-noshorthands,es-tabla]{babel}
\usepackage{lmodern}
\usepackage[margin=2.4cm]{geometry}
\usepackage{amsmath,amssymb,amsthm}
\usepackage{booktabs}
\usepackage{xcolor}
\usepackage{tikz}
\usepackage{algorithm}
\usepackage{algpseudocode}
\usepackage{enumitem}
\usepackage{float}
\usepackage[hidelinks]{hyperref}

\theoremstyle{definition}
\newtheorem{definicion}{Definición}

\floatname{algorithm}{Algoritmo}
\renewcommand{\algorithmicrequire}{\textbf{Entrada:}}
\renewcommand{\algorithmicensure}{\textbf{Salida:}}
\renewcommand{\listalgorithmname}{Lista de algoritmos}

% Etiquetas para distinguir el origen de cada afirmación
\newcommand{\req}{\textcolor{blue!60!black}{\textbf{[Requisito]}}}
\newcommand{\prop}{\textcolor{green!40!black}{\textbf{[Propuesta]}}}
\newcommand{\duda}{\textcolor{red!60!black}{\textbf{[Duda]}}}

\title{Representación del espacio de búsqueda y\\
       Optimización por Colonia de Hormigas (ACO)\\
       \large aplicadas al problema del ciclo hamiltoniano}
\author{Osvaldo Fabian Ramiro Balboa}
\date{Noviembre de 2026}

\begin{document}
\maketitle

\begin{abstract}
Este documento reúne la investigación previa (fase 1 del aula invertida) sobre dos temas:
cómo representar el grafo y las rutas candidatas sin producir soluciones ilegales, y cómo
funciona la Optimización por Colonia de Hormigas (ACO). Se incluye un ejemplo didáctico de
siete vértices, el pseudocódigo del método adaptado al ciclo hamiltoniano, las decisiones de
diseño posibles con sus dificultades y las fuentes que respaldan cada afirmación. Los números
del ejemplo son ilustrativos y \emph{no} son resultados experimentales sobre el FHCP.
\end{abstract}

\paragraph{Cómo leer este documento.} Para separar lo que pide la práctica de lo que
proponemos como equipo, algunas frases llevan una etiqueta: \req{} lo exige el enunciado,
\prop{} es una opción de diseño sugerida por las fuentes y \duda{} es una pregunta abierta
para resolver con el profesor o con el equipo. Las referencias entre corchetes, como [1],
remiten a la sección~\ref{sec:fuentes}.

\tableofcontents

% =====================================================================
\section{Términos básicos}
% =====================================================================

\begin{definicion}[Grafo]
Un grafo no dirigido es un par $G=(V,E)$, donde $V$ es un conjunto de $n$ \emph{vértices}
y $E$ es un conjunto de $m$ \emph{aristas}; cada arista $\{u,v\}$ une dos vértices distintos.
Dos vértices unidos por una arista son \emph{vecinos}; el \emph{grado} de un vértice es su
número de vecinos [6].
\end{definicion}

\begin{definicion}[Ciclo hamiltoniano]
Un ciclo hamiltoniano es una sucesión $(v_1,v_2,\dots,v_n)$ que contiene cada vértice de $V$
exactamente una vez, tal que $\{v_i,v_{i+1}\}\in E$ para $i=1,\dots,n-1$ y, además,
$\{v_n,v_1\}\in E$ (el \emph{cierre}). Tiene exactamente $n$ aristas. Decidir si un grafo
tiene uno es un problema NP-completo [7], por eso se usan metaheurísticas.
\end{definicion}

\begin{definicion}[Permutación]
Una permutación de $V$ es un ordenamiento de sus $n$ vértices en el que ninguno se repite
ni falta. Existen $n!$ permutaciones; si se fija el vértice inicial y se ignora el sentido
del recorrido, quedan $(n-1)!/2$ ciclos distintos posibles.
\end{definicion}

\begin{definicion}[Espacio de búsqueda y factibilidad]
El \emph{espacio de búsqueda} es el conjunto de todas las soluciones candidatas que el
algoritmo puede generar con la representación elegida. Una candidata es \emph{factible}
si cumple todas las reglas del problema; en este caso, si es un ciclo hamiltoniano.
\end{definicion}

% =====================================================================
\section{Tema general: representación del espacio de búsqueda}
% =====================================================================

La representación decide qué estructuras puede producir el algoritmo y, por lo tanto, qué
tan fácil es caer en soluciones ilegales. Larrañaga et al.~[4] muestran, para el problema
del viajero (TSP), que la misma metaheurística cambia mucho de comportamiento según la
codificación y los operadores. Hay que representar dos cosas distintas: el \emph{grafo}
(los datos del problema) y la \emph{ruta} (la solución candidata).

\subsection{Cómo guardar el grafo}

\paragraph{Matriz de adyacencia.} Es una tabla $A$ de $n\times n$ con $A[u][v]=1$ si la
arista $\{u,v\}$ existe y $0$ si no. Consultar si dos vértices están conectados cuesta
$O(1)$, pero la memoria es $O(n^2)$ aunque el grafo tenga pocas aristas [6].

\paragraph{Lista de adyacencia.} Para cada vértice se guarda la lista de sus vecinos.
Usa memoria $O(n+m)$ y permite recorrer los vecinos de un vértice en tiempo proporcional a
su grado, que es justo lo que hace una hormiga al elegir el siguiente paso. Saber si existe
una arista concreta cuesta $O(\text{grado})$ [6].

\begin{table}[h]
\centering
\caption{Comparación de representaciones del grafo.}
\label{tab:repr-grafo}
\begin{tabular}{@{}lll@{}}
\toprule
 & Matriz de adyacencia & Lista de adyacencia \\
\midrule
Memoria                         & $O(n^2)$ & $O(n+m)$ \\
¿Existe $\{u,v\}$?              & $O(1)$   & $O(\text{grado}(u))$ \\
Recorrer vecinos de $u$         & $O(n)$   & $O(\text{grado}(u))$ \\
Adecuada para                   & grafos densos o pequeños & grafos dispersos \\
\bottomrule
\end{tabular}
\end{table}

\paragraph{Dimensión real del FHCP.} Las 1001 instancias del FHCP tienen entre 66 y 9528
vértices [8,~9]. Para $n=9528$, una matriz tiene $9528^2\approx 9.08\times10^7$ celdas:
unos 91\,MB si se usa un byte por celda (unos 11\,MB si se usa un bit). Es manejable, pero
desperdicia memoria si el grafo es disperso.

\prop{} Usar \textbf{listas de adyacencia} para recorrer vecinos y, adicionalmente, un
\textbf{conjunto de aristas} (tabla hash de pares ordenados $(\min,\max)$) para verificar en
$O(1)$ el cierre y la validez de una ruta. \duda{} Medir el grado medio de las instancias
al descargarlas para confirmar que son dispersas.

\paragraph{Formato de los archivos.} El FHCP usa el formato HCP [9]: una cabecera, la lista
de aristas (una por línea) y un cierre. Las soluciones usan el mismo estilo con
\texttt{TYPE : TOUR} y una \texttt{TOUR\_SECTION} con los vértices en el orden del ciclo,
también terminada en \texttt{-1}. Esquema (los valores son ilustrativos):
\begin{center}
\begin{minipage}{0.42\textwidth}
\small
\begin{verbatim}
NAME : graph1
COMMENT : ...
TYPE : HCP
DIMENSION : 66
EDGE_DATA_FORMAT : EDGE_LIST
EDGE_DATA_SECTION
1 2
...
-1
EOF
\end{verbatim}
\end{minipage}
\end{center}
\duda{} Confirmar con un archivo real si los vértices se numeran desde 1.

\subsection{Cómo representar una ruta}

\req{} La práctica pide estructuras que representen permutaciones de los vértices sin
generar «cromosomas» ilegales, y que cada programa devuelva un arreglo de vértices como
solución propuesta.

\prop{} La representación más directa es la \textbf{permutación de vértices}
(\emph{path representation} en [4]): un arreglo de longitud $n$ como
$[1,7,6,5,4,3,2]$, que se lee como el recorrido $1\to7\to6\to5\to4\to3\to2\to1$.
Garantiza por construcción que cada vértice aparece una sola vez.

Otras codificaciones descritas en [4] y útiles para el equipo:
\begin{itemize}[nosep]
  \item \textbf{Adyacencia}: la posición $i$ guarda el sucesor del vértice $i$.
  \item \textbf{Ordinal}: cada gen indica una posición en una lista de vértices restantes;
        cualquier cruza de un punto produce una ruta válida.
  \item \textbf{Llaves aleatorias} (\emph{random keys}) [11]: un vector de números reales que
        se ordena para obtener la permutación. Es la puerta de entrada natural para
        algoritmos continuos como DE y PSO (integrantes 4 y 5).
\end{itemize}

\subsection{Por qué una permutación válida puede no ser un ciclo}

Una permutación cumple dos de las cuatro reglas del ciclo hamiltoniano: \emph{están todos
los vértices} y \emph{ninguno se repite}. Pero no garantiza las otras dos:
\begin{enumerate}[nosep]
  \item que cada par de vértices consecutivos esté unido por una arista del grafo, y
  \item que exista la arista de \emph{cierre} entre el último y el primero.
\end{enumerate}
En el TSP clásico el grafo es completo y toda permutación es una ruta; en el ciclo
hamiltoniano el grafo es incompleto (y en el FHCP, deliberadamente engañoso [8]), así que la
mayoría de las permutaciones usa aristas inexistentes.

\subsection{Cómo tratar rutas que no cumplen las reglas}

Coello Coello [5] clasifica las técnicas de manejo de restricciones; las tres que aplican aquí:

\begin{enumerate}
  \item \textbf{Evitar} (representación y operadores que preservan factibilidad): la ruta
        solo avanza por aristas existentes. ACO lo hace de forma natural, porque cada
        hormiga elige entre los \emph{vecinos} del vértice actual.
  \item \textbf{Penalizar}: se aceptan permutaciones arbitrarias y se mide cuánto les falta.
        \prop{} Función de costo
        \[
          f(\pi)=\sum_{i=1}^{n}\bigl[\{\pi_i,\pi_{(i \bmod n)+1}\}\notin E\bigr],
        \]
        es decir, el número de aristas faltantes incluido el cierre. Vale $0$ si y solo si
        $\pi$ es un ciclo hamiltoniano, y sirve como objetivo común para comparar GA, SA, DE
        y PSO.
  \item \textbf{Reparar}: transformar una candidata ilegal en una mejor. Por ejemplo, la
        \emph{rotación de Pósa} [10] toma un camino $v_1\dots v_k$ cuyo extremo $v_k$ es vecino
        de algún $v_i$, y lo reescribe como $v_1\dots v_i\,v_k v_{k-1}\dots v_{i+1}$, que tiene
        un extremo nuevo desde el cual intentar extender o cerrar.
\end{enumerate}

\paragraph{Comprobación de una ruta} (la misma para todos los algoritmos). Dada
$\pi=(\pi_1,\dots,\pi_n)$: (a) su longitud es $n$; (b) no hay repetidos (un conjunto de
tamaño $n$); (c) cada par consecutivo es arista; (d) $\{\pi_n,\pi_1\}$ es arista. Cuesta
$O(n)$ con el conjunto de aristas.

% =====================================================================
\section{Algoritmo: Optimización por Colonia de Hormigas (ACO)}
% =====================================================================

\subsection{Idea}

Las hormigas reales depositan feromona al caminar; los caminos cortos se recorren más
veces por unidad de tiempo, acumulan más feromona y atraen a más hormigas. ACO imita este
mecanismo: un conjunto de agentes construye soluciones paso a paso y se comunica de forma
indirecta (\emph{estigmergia}) a través de valores numéricos asociados a las aristas
[1,~3]. La primera versión fue el \emph{Ant System} (AS) [1]; luego aparecieron variantes
como \emph{Ant Colony System} (ACS) [2] y \emph{MAX--MIN Ant System} (MMAS) [3,~13].

\subsection{Construcción de rutas sin repetir vértices}

Cada hormiga $k$ empieza en un vértice y mantiene una \emph{memoria} (lista tabú) de los
vértices ya visitados. Estando en $i$, sus opciones son los vecinos aún no visitados
$N_i^k$. En AS elige $j\in N_i^k$ con probabilidad [1,~3]
\[
  p_{ij}^{k}=\frac{\tau_{ij}^{\alpha}\,\eta_{ij}^{\beta}}
                  {\sum_{l\in N_i^k}\tau_{il}^{\alpha}\,\eta_{il}^{\beta}},
\]
donde $\tau_{ij}$ es la feromona de la arista, $\eta_{ij}$ es la información heurística,
y $\alpha,\beta$ ponderan una contra otra. Como solo se consideran vecinos no visitados,
\textbf{la hormiga nunca repite vértices ni usa aristas inexistentes}: la representación
evita por construcción las soluciones ilegales.

\subsection{Feromonas: refuerzo y evaporación}

La feromona $\tau_{ij}$ es la memoria colectiva: cuánto ha convenido usar la arista
$\{i,j\}$ en el pasado. Tras cada iteración [1,~3]:
\[
  \tau_{ij}\leftarrow(1-\rho)\,\tau_{ij}+\sum_{k}\Delta\tau_{ij}^{k}.
\]
\begin{itemize}[nosep]
  \item \textbf{Evaporación} ($0<\rho\le1$): todas las aristas pierden una fracción de su
        feromona. Evita que la colonia se quede para siempre con las primeras rutas
        encontradas.
  \item \textbf{Depósito} ($\Delta\tau_{ij}^k$): las hormigas refuerzan las aristas que
        usaron, en proporción a la calidad de su ruta. En el TSP, $\Delta\tau=Q/L_k$
        ($L_k$ = longitud del recorrido).
\end{itemize}
Variantes: en ACS solo la mejor ruta deposita y además hay una actualización local que
reduce la feromona de las aristas recién usadas para diversificar [2]; en MMAS solo
deposita la mejor hormiga y $\tau$ se limita a $[\tau_{\min},\tau_{\max}]$ para evitar el
estancamiento [13].

\subsection{Información del grafo para elegir el siguiente paso}

En el TSP la heurística es $\eta_{ij}=1/d_{ij}$ (preferir ciudades cercanas). En el ciclo
hamiltoniano no hay distancias, así que hay que usar la \emph{estructura} del grafo.
\prop{} Una opción natural es la regla de Warnsdorff (preferir el vecino con menos vecinos
libres):
\[
  \eta_{ij}=\frac{1}{\max\bigl(1,\;|N(j)\setminus\text{visitados}|\bigr)}.
\]
La intuición: un vértice con pocos vecinos libres corre peligro de quedar aislado, así que
conviene visitarlo pronto. Otras señales útiles: si un vértice tiene grado 2, sus dos
aristas son \emph{obligatorias} en cualquier ciclo hamiltoniano, y conviene evitar los
movimientos que dejen algún vértice libre sin vecinos libres. \duda{} La regla de Warnsdorff proviene del recorrido del
caballo; no encontramos todavía una fuente que la evalúe con ACO sobre el FHCP.

\subsection{Callejones sin salida y ciclos que no cierran}

Al trabajar sobre un grafo incompleto aparecen dos fallas que no existen en el TSP:
\begin{enumerate}
  \item \textbf{Callejón sin salida}: la hormiga llega a un vértice cuyos vecinos ya fueron
        visitados, pero todavía faltan vértices. Su ruta es un camino de $k<n$ vértices.
  \item \textbf{Camino hamiltoniano sin cierre}: visita los $n$ vértices, pero el último no
        es vecino del primero.
\end{enumerate}
Opciones para tratarlas:
\begin{itemize}[nosep]
  \item \prop{} \emph{Puntuar} la ruta parcial con el número de aristas válidas recorridas
        (incluido el cierre si existe) dividido entre $n$, y depositar feromona en
        proporción a ese valor. Así una hormiga casi exitosa aporta información útil.
  \item \prop{} \emph{Reparar} con rotaciones de Pósa [10] antes de evaluar: si el extremo
        tiene un vecino ya visitado, rotar el camino y seguir extendiendo o intentar cerrar.
  \item \prop{} \emph{Retroceso limitado} (\emph{backtracking}): deshacer los últimos pasos
        y probar otro vecino, con un tope de intentos para no convertirlo en fuerza bruta.
  \item Wagner y Bruckstein [12] propusieron \emph{Hamiltonian(t)}, un proceso de una sola
        hormiga que marca vértices y elige según las marcas vecinas; muestran que los
        ciclos hamiltonianos son ciclos límite de ese proceso. Es un antecedente directo de
        hormigas aplicadas a este problema.
\end{itemize}

% =====================================================================
\section{Ejemplo didáctico con siete vértices}
% =====================================================================

\begin{figure}[h]
\centering
\begin{tikzpicture}[
  every node/.style={circle,draw,thick,minimum size=7mm,font=\small},
  ciclo/.style={very thick,blue!60!black},
  cuerda/.style={gray,dashed}]
  \foreach \v [count=\i from 0] in {1,...,7}
    \node (\v) at ({90-\i*360/7}:2.4) {\v};
  \draw[ciclo] (1)--(2)--(3)--(4)--(5)--(6)--(7)--(1);
  \draw[cuerda] (1)--(4) (2)--(6) (3)--(6);
\end{tikzpicture}
\caption{Grafo de ejemplo con $n=7$ y $m=10$. En azul, el único ciclo hamiltoniano; en
gris discontinuo, las tres aristas «trampa» $\{1,4\}$, $\{2,6\}$ y $\{3,6\}$.}
\label{fig:ejemplo}
\end{figure}

\paragraph{Representación.} Lista de adyacencia:
$1\!:\!\{2,4,7\}$, $2\!:\!\{1,3,6\}$, $3\!:\!\{2,4,6\}$, $4\!:\!\{1,3,5\}$, $5\!:\!\{4,6\}$,
$6\!:\!\{2,3,5,7\}$, $7\!:\!\{1,6\}$.

\paragraph{Razonamiento estructural.} Los vértices 5 y 7 tienen grado 2, así que las aristas
$\{4,5\}$, $\{5,6\}$, $\{6,7\}$ y $\{7,1\}$ son obligatorias. Entonces el vértice 6 ya usa dos aristas
(hacia 5 y 7) y quedan descartadas $\{2,6\}$ y $\{3,6\}$. Falta unir 1 con 4 pasando por 2 y
3: solo es posible como $1\!-\!2\!-\!3\!-\!4$. El ciclo hamiltoniano es único:
$1\!-\!2\!-\!3\!-\!4\!-\!5\!-\!6\!-\!7\!-\!1$ (lo confirmamos enumerando las $6!$
permutaciones que empiezan en 1: solo esta y su recorrido inverso tienen costo 0). Este es
el mismo tipo de estructura que vuelve difíciles a las instancias del FHCP: pocas aristas
correctas escondidas entre muchas tentadoras [8].

\paragraph{Costo de permutaciones (función de penalización).}
\begin{center}
\begin{tabular}{@{}lcl@{}}
\toprule
Permutación & Aristas faltantes & Motivo \\
\midrule
$[1,4,3,6,7,2,5]$ & 3 & faltan $\{7,2\}$, $\{2,5\}$ y el cierre $\{5,1\}$ \\
$[2,3,4,1,7,6,5]$ & 1 & todas existen salvo el cierre $\{5,2\}$ \\
$[1,2,3,4,5,6,7]$ & 0 & ciclo hamiltoniano \\
\bottomrule
\end{tabular}
\end{center}

\paragraph{Primer paso de una hormiga.} Parámetros $\alpha=1$, $\beta=2$, $\tau_0=1$ en
todas las aristas y la heurística de Warnsdorff. Desde el vértice 1, los vecinos libres de
2, 4 y 7 son 2, 2 y 1, respectivamente, así que $\eta=(\tfrac12,\tfrac12,1)$ y
\[
  p_{12}=\frac{0.25}{0.25+0.25+1}=0.167,\quad p_{14}=0.167,\quad p_{17}=0.667.
\]
La heurística empuja hacia el vértice 7, de grado 2, que es correcto. Desde 7 el único
vecino libre es 6. En 6 (visitados $\{1,7,6\}$) los candidatos son 2, 3 y 5 con
$\eta=(1,\tfrac12,1)$, de modo que $p_{62}=0.444$, $p_{63}=0.111$, $p_{65}=0.444$.

\paragraph{Tres hormigas de la primera iteración.}
\begin{itemize}[nosep]
  \item \textbf{Hormiga A}: $1\!-\!7\!-\!6\!-\!2\!-\!3\!-\!4\!-\!5$. Visita los 7 vértices,
        pero $\{5,1\}$ no existe: \emph{camino hamiltoniano sin cierre}. 6 aristas válidas.
  \item \textbf{Hormiga B}: $1\!-\!4\!-\!3\!-\!6\!-\!7$. En 7 sus vecinos (1 y 6) ya fueron
        visitados y faltan 2 y 5: \emph{callejón sin salida}. 4 aristas válidas.
  \item \textbf{Hormiga C}: $1\!-\!7\!-\!6\!-\!5\!-\!4\!-\!3\!-\!2\!-\!1$. Ciclo
        hamiltoniano: 7 aristas válidas y cierre correcto.
\end{itemize}
\textbf{Comprobación de C}: longitud 7, sin repetidos, las aristas
$\{1,7\},\{7,6\},\{6,5\},\{5,4\},\{4,3\},\{3,2\}$ existen y el cierre $\{2,1\}$ también.

\paragraph{Actualización de feromonas} con $\rho=0.5$ y depósito
$\Delta\tau=(\text{aristas válidas})/n$: A deposita $6/7=0.857$, B $4/7=0.571$ y C $1$.

\begin{center}
\begin{tabular}{@{}lcccccccccc@{}}
\toprule
Arista & 1-2 & 1-4 & 1-7 & 2-3 & 2-6 & 3-4 & 3-6 & 4-5 & 5-6 & 6-7 \\
\midrule
$\tau$ nueva & 1.500 & 1.071 & 2.357 & 2.357 & 1.357 & 2.929 & 1.071 & 2.357 & 1.500 & 2.929 \\
\bottomrule
\end{tabular}
\end{center}

En la segunda iteración, desde 6 (tras $1\!-\!7$) las probabilidades pasan a
$p_{62}=0.434$, $p_{63}=0.086$, $p_{65}=0.480$: el aprendizaje favorece la arista correcta,
pero lentamente, porque la hormiga A también reforzó la trampa $\{2,6\}$. Si se usa un
esquema elitista como MMAS, en el que solo deposita la mejor ruta (C), entonces
$\tau_{62}=0.5$, $\tau_{63}=0.5$ y $\tau_{65}=1.5$, y las probabilidades quedan en
$0.235$, $0.059$ y $0.706$. Lección: \textbf{en grafos engañosos importa mucho quién
deposita feromona}.

% =====================================================================
\section{Pseudocódigo y ajustes principales}
% =====================================================================

\begin{algorithm}[H]
\caption{ACO para el ciclo hamiltoniano (variante estilo MMAS)}
\label{alg:aco}
\begin{algorithmic}[1]
\Require Grafo $G=(V,E)$ como listas de adyacencia; parámetros $m,\alpha,\beta,\rho,
         \tau_{\min},\tau_{\max}$, límite de iteraciones $T$
\Ensure Arreglo de $n$ vértices que forma un ciclo hamiltoniano, o la mejor ruta parcial
\State $\tau_{ij}\gets\tau_{\max}$ para toda $\{i,j\}\in E$
\State $\mathit{mejor}\gets$ ruta vacía
\For{$t=1$ \textbf{to} $T$}
  \For{$k=1$ \textbf{to} $m$}
    \State elegir vértice inicial $s$ al azar; $R_k\gets[s]$; $\mathit{visitados}\gets\{s\}$
    \While{$|R_k|<n$}
      \State $i\gets$ último vértice de $R_k$
      \State $N\gets$ vecinos de $i$ no visitados
      \If{$N=\emptyset$}
        \State intentar reparar con rotación de Pósa; \textbf{si} no se puede, \textbf{break}
      \Else
        \State elegir $j\in N$ con probabilidad $\propto\tau_{ij}^{\alpha}\eta_{ij}^{\beta}$
        \State añadir $j$ a $R_k$ y a $\mathit{visitados}$
      \EndIf
    \EndWhile
    \State $q_k\gets$ (aristas válidas de $R_k$, incluido el cierre si existe)$/n$
    \If{$q_k=1$} \Return $R_k$ \Comment{ciclo hamiltoniano encontrado}
    \EndIf
    \If{$q_k>$ calidad de $\mathit{mejor}$} $\mathit{mejor}\gets R_k$
    \EndIf
  \EndFor
  \State $\tau_{ij}\gets(1-\rho)\tau_{ij}$ para toda arista \Comment{evaporación}
  \State $\tau_{ij}\gets\tau_{ij}+q_{\mathit{mejor}}$ para cada arista de $\mathit{mejor}$
         \Comment{depósito}
  \State acotar $\tau_{ij}$ al intervalo $[\tau_{\min},\tau_{\max}]$
\EndFor
\State \Return $\mathit{mejor}$
\end{algorithmic}
\end{algorithm}

\begin{table}[H]
\centering
\caption{Ajustes principales y valores de partida. Los valores de referencia son los que
Dorigo y Stützle [3] recomiendan \emph{para el TSP}; hay que calibrarlos para el FHCP.}
\label{tab:params}
\begin{tabular}{@{}p{2.7cm}p{6.9cm}p{4.4cm}@{}}
\toprule
Parámetro & Efecto & Referencia (TSP) \\
\midrule
$m$ (hormigas) & Más hormigas exploran más por iteración, pero cuestan más tiempo.
               & $m=n$ (AS, MMAS); $m=10$ (ACS) \\
$\alpha$ & Peso de la feromona (lo aprendido). Muy alto: convergencia prematura. & $1$ \\
$\beta$  & Peso de la heurística (la estructura del grafo). & $2$ a $5$ \\
$\rho$   & Evaporación. Alta: olvida rápido y explora; baja: explota lo aprendido.
         & $0.5$ (AS), $0.1$ (ACS), $0.02$ (MMAS) \\
$\tau_{\min},\tau_{\max}$ & Evitan que una arista quede con probabilidad casi 0 o casi 1. & MMAS \\
$q_0$    & (ACS) probabilidad de elegir directamente el mejor vecino. & $0.9$ \\
$T$      & Límite de iteraciones o de tiempo; necesario para la comparación. & ---\\
\bottomrule
\end{tabular}
\end{table}

% =====================================================================
\section{Aplicación al ciclo hamiltoniano: opciones, dificultades y dudas}
% =====================================================================

\subsection{Opciones de diseño}
\begin{enumerate}
  \item \prop{} \textbf{Grafo de construcción = grafo del problema}: las hormigas solo se
        mueven por aristas reales, así que nunca producen rutas ilegales; las fallas se
        reducen a «no llegó a todos» o «no cerró».
  \item \prop{} \textbf{Calidad} = fracción de aristas válidas (Sección~3.5). Es equivalente a
        $1-f(\pi)/n$ con la función de penalización de la Sección~2.4, lo que permite
        comparar ACO con GA, SA, DE y PSO con la misma métrica.
  \item \prop{} \textbf{Variante MMAS o ACS}: el ejemplo muestra que el depósito de todas las
        hormigas refuerza trampas; un esquema elitista con límites de feromona es más
        apropiado.
  \item \prop{} \textbf{Preprocesamiento}: fijar las aristas obligatorias de los vértices de
        grado 2 antes de lanzar la colonia.
  \item \prop{} \textbf{Búsqueda local}: combinar ACO con rotaciones de Pósa o 2-opt
        restringido a aristas existentes, como se hace con ACO+búsqueda local en el TSP [3].
\end{enumerate}

\subsection{Dificultades previstas}
\begin{itemize}[nosep]
  \item Las instancias del FHCP están construidas para engañar a heurísticas estándar [8];
        muchas aristas «atractivas» no pertenecen a ningún ciclo.
  \item Si el número de hormigas es $m=n$ y el grafo tiene miles de vértices, una iteración
        cuesta del orden de $n^2$ pasos; para $n\approx9500$ son unos $9\times10^7$ pasos.
  \item La feromona mide aristas, pero el problema es global: una arista puede estar en
        rutas casi completas y aun así no pertenecer a ningún ciclo hamiltoniano.
  \item Las rutas que fallan muy pronto dan poca información; la retroalimentación es escasa.
  \item Como referencia de dificultad, el ganador del concurso FHCP resolvió 985 de las
        1001 instancias [9]; no se espera que una metaheurística simple las resuelva todas.
\end{itemize}

\subsection{Dudas pendientes}
\begin{enumerate}[nosep]
  \item \duda{} ¿El profesor espera trabajar con las 1001 instancias o con un subconjunto
        (por ejemplo, las más pequeñas)?
  \item \duda{} ¿Qué límite de tiempo o de iteraciones usaremos para que la comparación sea
        justa entre los seis algoritmos?
  \item \duda{} ¿Se permite combinar ACO con búsqueda local o reparación, o se pide la
        versión «pura» del algoritmo?
  \item \duda{} ¿Cómo aplicar precisión, \emph{recall} y F1 (tema del integrante 4) a un
        problema donde todas las instancias son hamiltonianas? Una propuesta es tratar
        «encontró el ciclo» como predicción positiva y comparar contra la solución oficial.
  \item \duda{} ¿Acordamos en el equipo una única función de costo (aristas faltantes) para
        todos los algoritmos?
\end{enumerate}

% =====================================================================
\section{Fuentes consultadas}
\label{sec:fuentes}
% =====================================================================

\begin{enumerate}[label={[\arabic*]},leftmargin=*]
  \item M.~Dorigo, V.~Maniezzo y A.~Colorni. «Ant System: Optimization by a colony of
        cooperating agents». \emph{IEEE Trans. Systems, Man, and Cybernetics, Part B},
        26(1):29--41, 1996.
        \url{https://doi.org/10.1109/3477.484436}\\
        \emph{Respalda}: regla de transición probabilística, evaporación y depósito del
        Ant System.
  \item M.~Dorigo y L.~M.~Gambardella. «Ant Colony System: A cooperative learning approach to
        the traveling salesman problem». \emph{IEEE Trans. Evolutionary Computation},
        1(1):53--66, 1997.
        \url{https://doi.org/10.1109/4235.585892}\\
        \emph{Respalda}: variante ACS, actualización local y parámetro $q_0$.
  \item M.~Dorigo y T.~Stützle. \emph{Ant Colony Optimization}. MIT Press, 2004.
        \url{https://mitpress.mit.edu/9780262042192/ant-colony-optimization/}\\
        \emph{Respalda}: visión general de ACO, estigmergia, MMAS, valores de parámetros
        recomendados para el TSP y combinación con búsqueda local.
  \item P.~Larrañaga, C.~M.~H.~Kuijpers, R.~H.~Murga, I.~Inza y S.~Dizdarevic. «Genetic
        algorithms for the travelling salesman problem: A review of representations and
        operators». \emph{Artificial Intelligence Review}, 13(2):129--170, 1999.
        \url{https://doi.org/10.1023/A:1006529012972}\\
        \emph{Respalda}: representaciones de rutas (camino, adyacencia, ordinal) y su
        efecto en la búsqueda.
  \item C.~A.~Coello Coello. «Theoretical and numerical constraint-handling techniques used
        with evolutionary algorithms: a survey of the state of the art».
        \emph{Computer Methods in Applied Mechanics and Engineering}, 191(11--12):1245--1287,
        2002. \url{https://doi.org/10.1016/S0045-7825(01)00323-1}\\
        \emph{Respalda}: clasificación de técnicas para manejar restricciones (penalización,
        reparación, operadores que preservan factibilidad).
  \item T.~H.~Cormen, C.~E.~Leiserson, R.~L.~Rivest y C.~Stein. \emph{Introduction to
        Algorithms}, 3.ª ed., MIT Press, 2009, sección 22.1.
        ISBN 978-0-262-03384-8.\\
        \emph{Respalda}: definiciones de grafo y representaciones por matriz y por listas
        de adyacencia, con sus costos.
  \item R.~M.~Karp. «Reducibility among combinatorial problems». En \emph{Complexity of
        Computer Computations}, pp.~85--103, Plenum Press, 1972.
        \url{https://doi.org/10.1007/978-1-4684-2001-2_9}\\
        \emph{Respalda}: el ciclo hamiltoniano es NP-completo.
  \item M.~Haythorpe. «FHCP Challenge Set: The first set of structurally difficult instances
        of the Hamiltonian cycle problem». \emph{Bulletin of the ICA}, 83:98--107, 2018.
        \url{https://arxiv.org/abs/1902.10352}\\
        \emph{Respalda}: 1001 instancias hamiltonianas cuya dificultad proviene de la
        estructura, no solo del tamaño.
  \item Flinders Hamiltonian Cycle Project. «FHCP Challenge Set» (sitio oficial).
        \url{https://sites.flinders.edu.au/flinders-hamiltonian-cycle-project/fhcp-challenge-set/}\\
        \emph{Respalda}: rango de 66 a 9528 vértices, formato HCP y TOUR, resultado del
        concurso (985 de 1001 grafos resueltos por el equipo ganador).
  \item L.~Pósa. «Hamiltonian circuits in random graphs». \emph{Discrete Mathematics},
        14(4):359--364, 1976. \url{https://doi.org/10.1016/0012-365X(76)90068-6}\\
        \emph{Respalda}: técnica de rotación para extender caminos hacia ciclos hamiltonianos.
  \item J.~C.~Bean. «Genetic algorithms and random keys for sequencing and optimization».
        \emph{ORSA Journal on Computing}, 6(2):154--160, 1994.
        \url{https://doi.org/10.1287/ijoc.6.2.154}\\
        \emph{Respalda}: representación por llaves aleatorias.
  \item I.~A.~Wagner y A.~M.~Bruckstein. «Hamiltonian(t) -- An ant-inspired heuristic for
        recognizing Hamiltonian graphs». \emph{Proc. Congress on Evolutionary Computation
        (CEC 99)}, 1999. \url{https://doi.org/10.1109/CEC.1999.782656}\\
        \emph{Respalda}: antecedente de hormigas aplicadas específicamente al ciclo
        hamiltoniano.
  \item T.~Stützle y H.~H.~Hoos. «MAX--MIN Ant System». \emph{Future Generation Computer
        Systems}, 16(8):889--914, 2000.
        \url{https://doi.org/10.1016/S0167-739X(00)00043-1}\\
        \emph{Respalda}: depósito solo de la mejor hormiga y límites de feromona.
\end{enumerate}

\end{document}
